\section{Reinforcement Learning}
\setcounter{aufgabe}{0}

  % ===================================================================
% Aufgaben zu Reinforcement Learning
% Themen: Return und Diskontierung, Wertefunktionen,
%         Q-Learning (deterministisch und mit Lernrate),
%         nicht-deterministische Umgebungen, epsilon-greedy
% ===================================================================

% -------------------------------------------------------------------
% 1) Return und Diskontierung
% -------------------------------------------------------------------
\aufgabe{
Ein Agent durchläuft einen Spieldurchlauf und erhält ab Zeitpunkt $t$ die
Belohnungen
\[
  r_{t+1} = 2, \quad r_{t+2} = 0, \quad r_{t+3} = 4, \quad
  r_{t+4} = 10 ,
\]
danach ist der Durchlauf beendet. Der Return ist
\[
  G_t = \sum_{k=0}^{\infty} \gamma^k\, r_{t+k+1} .
\]

\begin{enumerate}[label=\alph*)]
  \item Berechnen Sie $G_t$ für $\gamma = 0$.
  \item Berechnen Sie $G_t$ für $\gamma = 0{,}5$.
  \item Berechnen Sie $G_t$ für $\gamma = 1$ und interpretieren Sie den
        Einfluss von $\gamma$.
\end{enumerate}
}

\loesung{
\begin{enumerate}[label=\alph*)]
  \item Für $\gamma = 0$ gilt $\gamma^0 = 1$ und $\gamma^k = 0$ für
  $k \geq 1$, also zählt nur die erste Belohnung:
  \[
    G_t = 1\cdot 2 = \textbf{2} .
  \]

  \item Für $\gamma = 0{,}5$:
  \[
    G_t = 2 + 0{,}5\cdot 0 + 0{,}25\cdot 4 + 0{,}125\cdot 10
        = 2 + 0 + 1 + 1{,}25 = \textbf{4{,}25} .
  \]

  \item Für $\gamma = 1$ werden alle Belohnungen gleich gewichtet:
  \[
    G_t = 2 + 0 + 4 + 10 = \textbf{16} .
  \]
  Kleines $\gamma$ macht den Agenten \textbf{kurzsichtig} (nur die nächste
  Belohnung zählt), großes $\gamma$ gewichtet \textbf{zukünftige}
  Belohnungen stark. Die späte Belohnung $10$ dominiert erst bei großem
  $\gamma$.
\end{enumerate}
}

% -------------------------------------------------------------------
% 2) Diskontierung unendlicher Belohnungen
% -------------------------------------------------------------------
\aufgabe{
Ein Agent erhält in \textbf{jedem} Zeitschritt konstant die Belohnung
$r = 3$, der Durchlauf ist unendlich lang. Der Return ist
\[
  G_t = \sum_{k=0}^{\infty} \gamma^k\, r , \qquad \gamma \in [0,1[ .
\]

\begin{enumerate}[label=\alph*)]
  \item Bestimmen Sie $G_t$ in geschlossener Form mithilfe der
        geometrischen Reihe.
  \item Berechnen Sie $G_t$ für $\gamma = 0{,}9$.
  \item Erklären Sie, warum die Summe für $\gamma = 1$ nicht mehr endlich
        ist.
\end{enumerate}
}

\loesung{
\begin{enumerate}[label=\alph*)]
  \item Mit der geometrischen Reihe
  $\sum_{k=0}^{\infty}\gamma^k = \tfrac{1}{1-\gamma}$ für $|\gamma| < 1$:
  \[
    G_t = r\sum_{k=0}^{\infty}\gamma^k = \frac{r}{1-\gamma} .
  \]

  \item Für $\gamma = 0{,}9$:
  \[
    G_t = \frac{3}{1-0{,}9} = \frac{3}{0{,}1} = \textbf{30} .
  \]

  \item Für $\gamma = 1$ gilt $\gamma^k = 1$ für alle $k$, also
  \[
    G_t = \sum_{k=0}^{\infty} 3 = \infty .
  \]
  Die Reihe divergiert. Die Diskontierung mit $\gamma < 1$ sichert also die
  \textbf{Konvergenz} des Returns bei unendlich langen Durchläufen.
\end{enumerate}
}

% -------------------------------------------------------------------
% 3) Q-Learning: deterministisch (ein Update)
% -------------------------------------------------------------------
\aufgabe{
In einer \textbf{deterministischen} Umgebung wird der Q-Wert nach
\[
  Q(s_t, a_t) = r_t + \gamma\,\max_a Q(s_{t+1}, a)
\]
aktualisiert. Der Agent wählt in $s_1$ die Aktion $a_1$, erhält
$r_t = 3$ und landet in $s_2$. Dort gilt
\[
  Q(s_2, a_1) = 8, \qquad Q(s_2, a_2) = 12, \qquad Q(s_2, a_3) = -4 .
\]
Der Diskontfaktor ist $\gamma = 0{,}9$.

\begin{enumerate}[label=\alph*)]
  \item Bestimmen Sie $\max_a Q(s_2, a)$.
  \item Berechnen Sie den neuen Wert $Q(s_1, a_1)$.
  \item Wie ändert sich das Ergebnis für $\gamma = 0{,}5$?
\end{enumerate}
}

\loesung{
\begin{enumerate}[label=\alph*)]
  \item Der maximale Folge-Q-Wert ist
  \[
    \max_a Q(s_2, a) = \max\{8, 12, -4\} = 12 .
  \]

  \item Mit $\gamma = 0{,}9$:
  \[
    Q(s_1, a_1) = 3 + 0{,}9\cdot 12 = 3 + 10{,}8 = \textbf{13{,}8} .
  \]

  \item Mit $\gamma = 0{,}5$:
  \[
    Q(s_1, a_1) = 3 + 0{,}5\cdot 12 = 3 + 6 = \textbf{9} .
  \]
  Der kleinere Diskontfaktor gewichtet den zukünftigen Erfolg schwächer,
  daher fällt der Q-Wert niedriger aus.
\end{enumerate}
}

% -------------------------------------------------------------------
% 4) Q-Learning: Update mit Lernrate (mehrere Runden)
% -------------------------------------------------------------------
\aufgabe{
In einer \textbf{nicht-deterministischen} Umgebung lautet die Update-Regel
\[
  Q(s_t, a_t) = (1-\alpha)\,Q(s_t, a_t)
    + \alpha\bigl(r_t + \gamma\,\max_a Q(s_{t+1}, a)\bigr)
\]
mit $\alpha = 0{,}3$ und $\gamma = 0{,}9$. Der Startwert ist
$Q(s_1, a_1) = 0$.

\begin{enumerate}[label=\alph*)]
  \item Runde 1: Die Aktion liefert
        $r_t + \gamma\max_a Q(s_{t+1}, a) = 20$. Berechnen Sie den neuen
        Q-Wert.
  \item Runde 2: Nun liefert dieselbe Aktion
        $r_t + \gamma\max_a Q(s_{t+1}, a) = -30$ (der Agent stirbt).
        Berechnen Sie den neuen Q-Wert.
  \item Erklären Sie, welche Rolle die Lernrate $\alpha$ bei zufälligen
        Ergebnissen spielt.
\end{enumerate}
}

\loesung{
\begin{enumerate}[label=\alph*)]
  \item Runde 1:
  \[
    Q(s_1, a_1) = 0{,}7\cdot 0 + 0{,}3\cdot 20 = \textbf{6} .
  \]

  \item Runde 2 mit dem neuen Startwert $6$:
  \[
    Q(s_1, a_1) = 0{,}7\cdot 6 + 0{,}3\cdot(-30)
                = 4{,}2 - 9 = \textbf{-4{,}8} .
  \]

  \item Die Lernrate $\alpha$ gewichtet die \textbf{aktuelle} Erfahrung
  gegen den bisherigen Wert. Zufällige Schwankungen werden über viele
  Runden \textbf{geglättet}, sodass sich $Q$ dem Erwartungswert nähert.
  Für $\alpha = 1$ würde jede Runde den alten Wert komplett überschreiben,
  was nur in \textbf{deterministischen} Umgebungen sinnvoll ist.
\end{enumerate}
}

% -------------------------------------------------------------------
% 5) Q-Learning: vollstaendiger Pfad
% -------------------------------------------------------------------
\aufgabe{
Betrachtet wird eine deterministische Kette von Zuständen
$s_1 \to s_2 \to s_3 \to s_4$ (Ziel). Die Belohnungen der Übergänge sind
\[
  r(s_1 \to s_2) = 0, \quad
  r(s_2 \to s_3) = 0, \quad
  r(s_3 \to s_4) = 10 .
\]
Im Ziel $s_4$ gilt $Q(s_4, \cdot) = 0$. Alle Q-Werte starten bei $0$,
$\gamma = 0{,}9$, und es wird deterministisch aktualisiert
($Q(s,a) = r + \gamma\max_{a'} Q(s', a')$). Es gibt pro Zustand nur eine
Aktion.

\begin{enumerate}[label=\alph*)]
  \item Aktualisieren Sie die Q-Werte in der Reihenfolge
        $s_3, s_2, s_1$ (rückwärts). Geben Sie die Werte an.
  \item Erklären Sie, warum ein einziger Vorwärts-Durchlauf nicht ausreicht,
        um alle Q-Werte auf ihren Endwert zu bringen.
\end{enumerate}
}

\loesung{
\begin{enumerate}[label=\alph*)]
  \item Rückwärts, beginnend beim zielnahen Zustand:
  \[
    Q(s_3, \cdot) = 10 + 0{,}9\cdot Q(s_4, \cdot) = 10 + 0 = 10 ,
  \]
  \[
    Q(s_2, \cdot) = 0 + 0{,}9\cdot Q(s_3, \cdot) = 0{,}9\cdot 10 = 9 ,
  \]
  \[
    Q(s_1, \cdot) = 0 + 0{,}9\cdot Q(s_2, \cdot) = 0{,}9\cdot 9 = 8{,}1 .
  \]

  \item In einem einzelnen Vorwärts-Durchlauf ist beim Update von $s_1$ der Wert
  $Q(s_2, \cdot)$ noch $0$, sodass $Q(s_1, \cdot)$ zunächst $0$ bleibt. Die
  Belohnung aus dem Ziel muss sich erst \textbf{Schritt für Schritt}
  rückwärts durch die Kette ausbreiten. Nach genügend vielen Durchläufen
  konvergieren die Werte gegen $10$, $9$ und $8{,}1$.
\end{enumerate}
}

% -------------------------------------------------------------------
% 6) Erwartungswert des Q-Updates
% -------------------------------------------------------------------
\aufgabe{
Eine Aktion $a$ im Zustand $s$ führt \textbf{zufällig} zu zwei Ergebnissen:
mit Wahrscheinlichkeit $0{,}8$ zum Zustand $s'$ mit Reward $r = 5$ und mit
Wahrscheinlichkeit $0{,}2$ zum Zustand $s''$ mit Reward $r = -10$. Es gelte
\[
  \max_a Q(s', a) = 20, \qquad \max_a Q(s'', a) = 0, \qquad \gamma = 0{,}9 .
\]

\begin{enumerate}[label=\alph*)]
  \item Berechnen Sie den Zielwert
        $y = r + \gamma\max_a Q(\cdot, a)$ für beide Ergebnisse.
  \item Berechnen Sie den Erwartungswert $E[y]$.
  \item Gegen welchen Wert konvergiert $Q(s,a)$, wenn das Q-Learning über
        viele Runden mittelt?
\end{enumerate}
}

\loesung{
\begin{enumerate}[label=\alph*)]
  \item Ergebnis $s'$:
  \[
    y_1 = 5 + 0{,}9\cdot 20 = 23 .
  \]
  Ergebnis $s''$:
  \[
    y_2 = -10 + 0{,}9\cdot 0 = -10 .
  \]

  \item Erwartungswert:
  \[
    E[y] = 0{,}8\cdot 23 + 0{,}2\cdot(-10) = 18{,}4 - 2 = \textbf{16{,}4} .
  \]

  \item Da die Update-Regel mit Lernrate über viele Runden den
  Erwartungswert approximiert, konvergiert
  \[
    Q(s,a) \to E[y] = 16{,}4 .
  \]
  Das einzelne zufällige Ergebnis wird also durch die Mittelung geglättet.
\end{enumerate}
}
